\chapter{Challenge 1: Spurverfolgung}
\label{ch:Challenge_1}
% Einleitung zur Challenge 1
In der ersten Aufgabe des Projekts bestand die Anforderung darin, den Duckiebot so zu programmieren, dass er sich innerhalb der markierten Fahrspuren in der Modellumgebung Duckietown autonom fortbewegt. Die Fahrspuren sind durch eine durchgezogene weiße Linie auf der rechten Seite und eine gestrichelte gelbe Linie auf der linken Seite definiert. Zusätzlich befinden sich an bestimmten Stellen rote Querstreifen, an denen der Duckiebot kurz anhalten und nach zwei Sekunden automatisch weiterfahren soll. Ziel dieser Aufgabe war die Entwicklung eines grundlegenden Spurführungsalgorithmus auf Basis visueller Daten, der es dem Fahrzeug ermöglicht, kontinuierlich seine Position relativ zur Fahrbahn zu bestimmen und entsprechend zu korrigieren. Diese Funktion bildet die Grundlage für sämtliche nachfolgenden Herausforderungen und ist essenziell für die autonome Navigation innerhalb des Gesamtsystems. 